% Auto-generated by statistical_analysis.py
% Comprehensive results: subgroup alignment + statistical analyses.
% Metrics: ModalMAE (lower better) and EAD (lower better).
% Question split: Perception = Q1, Q10 per DS; Information = Q2--Q9 per DS.
% Subgroup classification (humans): doctor/ER visit weekly|monthly = high,
%                                   yearly|never                  = low.

\section*{1. Subgroup-level Alignment with Human Responses}
Tables~\ref{tab:perc-distance-alignment}--\ref{tab:info-expdist-alignment} report distance-based alignment between each LLM subgroup and the matched human subgroup (e.g., LLM Edu-Low vs.\ Human Edu-Low). Modal MAE measures the gap between modal answers; Expected Absolute Distance (EAD) integrates over the full marginal distributions. Lower values indicate better alignment; the maximum possible distance is 4 under the ordinal coding A\,$=$\,0$\ldots$E\,$=$\,4.
Each table includes a \emph{Human (split-half)} baseline row: human respondents in each subgroup are randomly split into two halves and the same metric is computed between halves, averaged over $200$ bootstrap resamples. This baseline is the natural noise floor caused by within-subgroup human disagreement and finite sample size; any model whose row approaches this baseline is statistically indistinguishable from a second independent human sample of the same size.

\begin{table}[htbp]
\centering
\scriptsize
\setlength{\tabcolsep}{3.5pt}
\caption{Distance-based alignment for Perception-based questions (Q1 and Q10 per DS, $n{=}8$): Modal MAE between each LLM subgroup and the matched human subgroup. Lower is better.}
\label{tab:perc-distance-alignment}
\begin{tabular}{lcccccccc}
\toprule
\textbf{Model}
& \multicolumn{2}{c}{\textbf{Education}}
& \multicolumn{2}{c}{\textbf{Sex}}
& \multicolumn{2}{c}{\textbf{Outpatient}}
& \textbf{ED Visits}
& \textbf{Max} \\
\cmidrule(lr){2-3}\cmidrule(lr){4-5}\cmidrule(lr){6-7}
& Low & High & Male & Female & Low & High & Low & \\
\midrule
Sonnet 4.5 & 1.62 & 0.75 & 0.75 & 0.62 & 0.75 & 0.88 & 0.62 & 4 \\
Opus 4.6 & 1.38 & 1.25 & 1.12 & 1.12 & 1.00 & 0.38 & 1.00 & 4 \\
GPT-5.2 & 1.12 & 1.25 & 0.88 & 0.50 & 0.75 & 1.00 & 0.88 & 4 \\
GPT-4.1 & 1.38 & 1.12 & 1.25 & 0.62 & 0.88 & 1.12 & 1.12 & 4 \\
\midrule
\textit{Human (split-half)} & 0.39 & 0.64 & 0.27 & 0.59 & 0.47 & 0.82 & 0.36 & 4 \\
\bottomrule
\end{tabular}
\end{table}
\begin{table}[htbp]
\centering
\scriptsize
\setlength{\tabcolsep}{3.5pt}
\caption{Distribution-level alignment for Perception-based questions: expected absolute ordinal distance $\mathbb{E}[|\mathrm{ord}(\hat{Y}_{\mathrm{LLM}})-\mathrm{ord}(Y_{\mathrm{Human}})|]$ over matched subgroups.}
\label{tab:perc-expdist-alignment}
\begin{tabular}{lcccccccc}
\toprule
\textbf{Model}
& \multicolumn{2}{c}{\textbf{Education}}
& \multicolumn{2}{c}{\textbf{Sex}}
& \multicolumn{2}{c}{\textbf{Outpatient}}
& \textbf{ED Visits}
& \textbf{Max} \\
\cmidrule(lr){2-3}\cmidrule(lr){4-5}\cmidrule(lr){6-7}
& Low & High & Male & Female & Low & High & Low & \\
\midrule
Sonnet 4.5 & 1.30 & 0.85 & 0.97 & 1.02 & 0.98 & 1.08 & 0.97 & 4 \\
Opus 4.6 & 1.11 & 1.23 & 1.03 & 1.05 & 1.04 & 1.02 & 1.03 & 4 \\
GPT-5.2 & 1.03 & 1.08 & 0.92 & 0.98 & 0.96 & 0.93 & 0.95 & 4 \\
GPT-4.1 & 1.11 & 1.18 & 0.99 & 1.02 & 1.01 & 0.99 & 1.01 & 4 \\
\midrule
\textit{Human (split-half)} & 0.91 & 0.91 & 0.83 & 0.96 & 0.92 & 0.88 & 0.94 & 4 \\
\bottomrule
\end{tabular}
\end{table}
\begin{table}[htbp]
\centering
\scriptsize
\setlength{\tabcolsep}{3.5pt}
\caption{Distance-based alignment for Information-based questions (Q2--Q9 per DS, $n{=}32$): Modal MAE between each LLM subgroup and the matched human subgroup.}
\label{tab:info-distance-alignment}
\begin{tabular}{lcccccccc}
\toprule
\textbf{Model}
& \multicolumn{2}{c}{\textbf{Education}}
& \multicolumn{2}{c}{\textbf{Sex}}
& \multicolumn{2}{c}{\textbf{Outpatient}}
& \textbf{ED Visits}
& \textbf{Max} \\
\cmidrule(lr){2-3}\cmidrule(lr){4-5}\cmidrule(lr){6-7}
& Low & High & Male & Female & Low & High & Low & \\
\midrule
Sonnet 4.5 & 1.20 & 1.06 & 1.21 & 1.00 & 1.25 & 1.00 & 1.14 & 4 \\
Opus 4.6 & 0.72 & 0.68 & 0.68 & 0.57 & 0.86 & 0.57 & 0.69 & 4 \\
GPT-5.2 & 0.88 & 0.65 & 0.71 & 0.68 & 0.86 & 0.71 & 0.62 & 4 \\
GPT-4.1 & 1.00 & 0.90 & 0.75 & 0.71 & 0.79 & 0.79 & 0.76 & 4 \\
\midrule
\textit{Human (split-half)} & 0.55 & 0.32 & 0.52 & 0.37 & 0.40 & 0.41 & 0.35 & 4 \\
\bottomrule
\end{tabular}
\end{table}
\begin{table}[htbp]
\centering
\scriptsize
\setlength{\tabcolsep}{3.5pt}
\caption{Distribution-level alignment for Information-based questions: expected absolute ordinal distance over matched subgroups.}
\label{tab:info-expdist-alignment}
\begin{tabular}{lcccccccc}
\toprule
\textbf{Model}
& \multicolumn{2}{c}{\textbf{Education}}
& \multicolumn{2}{c}{\textbf{Sex}}
& \multicolumn{2}{c}{\textbf{Outpatient}}
& \textbf{ED Visits}
& \textbf{Max} \\
\cmidrule(lr){2-3}\cmidrule(lr){4-5}\cmidrule(lr){6-7}
& Low & High & Male & Female & Low & High & Low & \\
\midrule
Sonnet 4.5 & 0.95 & 1.01 & 1.05 & 1.00 & 1.11 & 0.98 & 1.07 & 4 \\
Opus 4.6 & 0.69 & 0.76 & 0.81 & 0.75 & 0.86 & 0.74 & 0.81 & 4 \\
GPT-5.2 & 0.80 & 0.78 & 0.86 & 0.78 & 0.89 & 0.82 & 0.86 & 4 \\
GPT-4.1 & 0.85 & 0.89 & 0.93 & 0.89 & 0.96 & 0.93 & 0.96 & 4 \\
\midrule
\textit{Human (split-half)} & 0.75 & 0.59 & 0.70 & 0.60 & 0.66 & 0.66 & 0.67 & 4 \\
\bottomrule
\end{tabular}
\end{table}
\section*{2. Stability Across $N{=}10$ Iterations}
Table~\ref{tab:stability} quantifies the noise reduction afforded by ten independent runs per model. Standard errors of the mean (SEM) are uniformly small (all $\le 0.017$, most $<0.005$), confirming that the ten-iteration average is a stable estimator of each model's expected behavior. EAD is approximately an order of magnitude more stable than Modal MAE, because the latter is a discrete plug-in statistic sensitive to mode flips.

\begin{table}[htbp]
\centering
\scriptsize
\setlength{\tabcolsep}{4pt}
\caption{Stability metrics over $N{=}10$ iterations per model. Lower is better for all columns. ModalMAE and EAD are averaged over the 8 matched subgroups.}
\label{tab:stability}
\begin{tabular}{llrrrr}
\toprule
\textbf{Metric} & \textbf{Type} & \textbf{Model} & \textbf{Mean} & \textbf{SD} & \textbf{SEM} \\
\midrule
ModalMAE & Perception & Sonnet 4.5 & 0.854 & 0.011 & 0.004 \\
 &  & Opus 4.6 & 1.032 & 0.008 & 0.002 \\
 &  & GPT-5.2 & 0.895 & 0.013 & 0.004 \\
 &  & GPT-4.1 & 0.973 & 0.052 & 0.016 \\
\cmidrule(lr){2-6}
 & Information & Sonnet 4.5 & 1.123 & 0.010 & 0.003 \\
 &  & Opus 4.6 & 0.678 & 0.006 & 0.002 \\
 &  & GPT-5.2 & 0.740 & 0.013 & 0.004 \\
 &  & GPT-4.1 & 0.810 & 0.012 & 0.004 \\
\cmidrule(lr){2-6}
EAD & Perception & Sonnet 4.5 & 1.024 & 0.004 & 0.001 \\
 &  & Opus 4.6 & 1.074 & 0.002 & 0.000 \\
 &  & GPT-5.2 & 0.979 & 0.003 & 0.001 \\
 &  & GPT-4.1 & 1.043 & 0.003 & 0.001 \\
\cmidrule(lr){2-6}
 & Information & Sonnet 4.5 & 1.024 & 0.003 & 0.001 \\
 &  & Opus 4.6 & 0.772 & 0.001 & 0.000 \\
 &  & GPT-5.2 & 0.826 & 0.006 & 0.002 \\
 &  & GPT-4.1 & 0.917 & 0.003 & 0.001 \\
\midrule
\bottomrule
\end{tabular}
\end{table}
\section*{3. Normality of Per-Iteration Score Distributions}
Table~\ref{tab:shapiro} reports Shapiro--Wilk results on the ten per-iteration scalar scores. EAD scores are normally distributed across all four models and both question types ($p>0.17$ throughout), so a parametric one-way ANOVA is appropriate. Modal MAE distributions occasionally violate normality on perception questions (Sonnet 4.5: $p\approx10^{-7}$; Opus 4.6: $p=0.004$); the corresponding ANOVA on Modal MAE/Perception should therefore be read alongside the K--S evidence.

\begin{table}[htbp]
\centering
\scriptsize
\setlength{\tabcolsep}{4pt}
\caption{Shapiro--Wilk normality test on the $N{=}10$ per-iteration scores for each model. $p > 0.05$ fails to reject normality (parametric tests appropriate).}
\label{tab:shapiro}
\begin{tabular}{llrrl}
\toprule
\textbf{Metric} & \textbf{Type} & \textbf{Model} & $W$ & $p$ \\
\midrule
ModalMAE & Perception & Sonnet 4.5 & 0.3657 & $1.00\times 10^{-7}$ \\
 &  & Opus 4.6 & 0.5093 & $4.67\times 10^{-6}$ \\
 &  & GPT-5.2 & 0.8325 & 0.0359 \\
 &  & GPT-4.1 & 0.9028 & 0.2354 \\
\cmidrule(lr){2-5}
 & Information & Sonnet 4.5 & 0.9111 & 0.2889 \\
 &  & Opus 4.6 & 0.9103 & 0.2833 \\
 &  & GPT-5.2 & 0.8960 & 0.1979 \\
 &  & GPT-4.1 & 0.8837 & 0.1440 \\
\cmidrule(lr){2-5}
EAD & Perception & Sonnet 4.5 & 0.9615 & 0.8030 \\
 &  & Opus 4.6 & 0.8924 & 0.1805 \\
 &  & GPT-5.2 & 0.9542 & 0.7187 \\
 &  & GPT-4.1 & 0.9490 & 0.6562 \\
\cmidrule(lr){2-5}
 & Information & Sonnet 4.5 & 0.9254 & 0.4040 \\
 &  & Opus 4.6 & 0.9845 & 0.9846 \\
 &  & GPT-5.2 & 0.9612 & 0.7997 \\
 &  & GPT-4.1 & 0.9113 & 0.2902 \\
\midrule
\bottomrule
\end{tabular}
\end{table}
\section*{4. Distributional Alignment with Humans (K--S)}
Table~\ref{tab:ks} reports two-sample Kolmogorov--Smirnov tests on the pooled ordinal LLM responses against the human responses, separately for perception and information questions. With sample sizes in the tens of thousands all $p$-values are highly significant; the supremum statistic $D$ is therefore the appropriate ranking quantity. For perception questions, GPT-5.2 achieves the closest match to humans ($D{=}0.070$), followed by GPT-4.1 ($D{=}0.093$); Sonnet 4.5 is farthest ($D{=}0.324$). For information questions Opus 4.6 is closest ($D{=}0.131$) and Sonnet 4.5 is again farthest ($D{=}0.221$).
The bottom row of each block reports the within-human \emph{Human (split-half)} K--S baseline (mean over $200$ random splits): $D\approx0.045$ for perception and $D\approx0.028$ for information, with mean $p\approx0.82$. This is the intrinsic floor due to finite human sample size. \textbf{No model reaches this floor}: even the best model (GPT-5.2 perception, $D{=}0.070$) is roughly $1.5\times$ the human noise level, indicating that meaningful headroom remains for closing the human-LLM gap.

\begin{table}[htbp]
\centering
\scriptsize
\setlength{\tabcolsep}{4pt}
\caption{Two-sample Kolmogorov--Smirnov tests comparing pooled LLM response distributions against human responses on ordinal encoding (A{=}0$\ldots$E{=}4). $D$ is the supremum CDF difference; small $p$ indicates distributional mismatch.}
\label{tab:ks}
\begin{tabular}{llrrrr}
\toprule
\textbf{Type} & \textbf{Model} & $n_{\text{LLM}}$ & $n_{\text{Hum}}$ & $D$ & $p$ \\
\midrule
Perception & Sonnet 4.5 & 15360 & 586 & 0.3239 & $3.61\times 10^{-53}$ \\
 & Opus 4.6 & 15360 & 586 & 0.1357 & $1.60\times 10^{-9}$ \\
 & GPT-5.2 & 15360 & 586 & 0.0699 & 0.0077 \\
 & GPT-4.1 & 15360 & 586 & 0.0925 & 0.0001 \\
 & \textit{Human (split-half)} & --- & --- & 0.0454 & 0.8172 \\
\midrule
Information & Sonnet 4.5 & 57880 & 1373 & 0.2207 & $7.57\times 10^{-58}$ \\
 & Opus 4.6 & 60077 & 1373 & 0.1310 & $1.54\times 10^{-20}$ \\
 & GPT-5.2 & 58981 & 1373 & 0.1624 & $2.19\times 10^{-31}$ \\
 & GPT-4.1 & 59152 & 1373 & 0.1428 & $2.47\times 10^{-24}$ \\
 & \textit{Human (split-half)} & --- & --- & 0.0278 & 0.8235 \\
\bottomrule
\end{tabular}
\end{table}
\section*{5. Inter-Model Benchmarking (ANOVA + Tukey HSD)}
Table~\ref{tab:anova} reports the omnibus one-way ANOVA over the four models (degrees of freedom $3,36$). All four metric/type combinations yield $F\ge65$ and $p\le10^{-14}$, confirming significant inter-model differences. Pairwise Tukey HSD tests are reported in Table~\ref{tab:tukey}. The post-hoc pattern reveals a clean dichotomy:
\begin{itemize}\itemsep0pt
\item \textbf{Information-based questions:} Opus 4.6 statistically dominates all other models on both Modal MAE and EAD ($p_{\text{adj}}<10^{-4}$ for every comparison), with GPT-5.2 second, GPT-4.1 third, and Sonnet 4.5 a clear last.
\item \textbf{Perception-based questions:} GPT-5.2 has the lowest EAD (significant against every competitor), while Sonnet 4.5 has the lowest Modal MAE. This split between metrics suggests Sonnet 4.5 nails the modal answer but allocates probability mass less faithfully than GPT-5.2.
\end{itemize}

\begin{table}[htbp]
\centering
\scriptsize
\setlength{\tabcolsep}{5pt}
\caption{One-way ANOVA across the four models (each with $N{=}10$ iteration scores). A significant $p$ indicates that at least one model differs from the others.}
\label{tab:anova}
\begin{tabular}{llrr}
\toprule
\textbf{Metric} & \textbf{Type} & $F(3,36)$ & $p$ \\
\midrule
ModalMAE & Perception & 83.044 & $3.05\times 10^{-16}$ \\
ModalMAE & Information & 3601.324 & $1.18\times 10^{-44}$ \\
EAD & Perception & 1600.066 & $2.40\times 10^{-38}$ \\
EAD & Information & 8166.188 & $4.85\times 10^{-51}$ \\
\bottomrule
\end{tabular}
\end{table}
\begin{table}[htbp]
\centering
\scriptsize
\setlength{\tabcolsep}{4pt}
\caption{Tukey HSD post-hoc pairwise comparisons between models (family-wise $\alpha{=}0.05$). Negative `meandiff' means Model\,1 has a lower (better) score than Model\,2. Reject indicates a statistically significant difference.}
\label{tab:tukey}
\begin{tabular}{llllrrl}
\toprule
\textbf{Metric} & \textbf{Type} & \textbf{Model 1} & \textbf{Model 2} & \textbf{meandiff} & $p_{\text{adj}}$ & \textbf{Reject} \\
\midrule
ModalMAE & Perception & GPT-4.1 & GPT-5.2 & -0.079 & 0.0010 & \textbf{Yes} \\
 &  & GPT-4.1 & Opus 4.6 & 0.059 & 0.0010 & \textbf{Yes} \\
 &  & GPT-4.1 & Sonnet 4.5 & -0.120 & 0.0010 & \textbf{Yes} \\
 &  & GPT-5.2 & Opus 4.6 & 0.138 & 0.0010 & \textbf{Yes} \\
 &  & GPT-5.2 & Sonnet 4.5 & -0.041 & 0.0108 & \textbf{Yes} \\
 &  & Opus 4.6 & Sonnet 4.5 & -0.179 & 0.0010 & \textbf{Yes} \\
\cmidrule(lr){2-7}
 & Information & GPT-4.1 & GPT-5.2 & -0.070 & 0.0010 & \textbf{Yes} \\
 &  & GPT-4.1 & Opus 4.6 & -0.133 & 0.0010 & \textbf{Yes} \\
 &  & GPT-4.1 & Sonnet 4.5 & 0.312 & 0.0010 & \textbf{Yes} \\
 &  & GPT-5.2 & Opus 4.6 & -0.063 & 0.0010 & \textbf{Yes} \\
 &  & GPT-5.2 & Sonnet 4.5 & 0.382 & 0.0010 & \textbf{Yes} \\
 &  & Opus 4.6 & Sonnet 4.5 & 0.445 & 0.0010 & \textbf{Yes} \\
\cmidrule(lr){2-7}
EAD & Perception & GPT-4.1 & GPT-5.2 & -0.064 & 0.0010 & \textbf{Yes} \\
 &  & GPT-4.1 & Opus 4.6 & 0.030 & 0.0010 & \textbf{Yes} \\
 &  & GPT-4.1 & Sonnet 4.5 & -0.019 & 0.0010 & \textbf{Yes} \\
 &  & GPT-5.2 & Opus 4.6 & 0.094 & 0.0010 & \textbf{Yes} \\
 &  & GPT-5.2 & Sonnet 4.5 & 0.045 & 0.0010 & \textbf{Yes} \\
 &  & Opus 4.6 & Sonnet 4.5 & -0.049 & 0.0010 & \textbf{Yes} \\
\cmidrule(lr){2-7}
 & Information & GPT-4.1 & GPT-5.2 & -0.090 & 0.0010 & \textbf{Yes} \\
 &  & GPT-4.1 & Opus 4.6 & -0.144 & 0.0010 & \textbf{Yes} \\
 &  & GPT-4.1 & Sonnet 4.5 & 0.107 & 0.0010 & \textbf{Yes} \\
 &  & GPT-5.2 & Opus 4.6 & -0.054 & 0.0010 & \textbf{Yes} \\
 &  & GPT-5.2 & Sonnet 4.5 & 0.197 & 0.0010 & \textbf{Yes} \\
 &  & Opus 4.6 & Sonnet 4.5 & 0.251 & 0.0010 & \textbf{Yes} \\
\midrule
\bottomrule
\end{tabular}
\end{table}
\section*{6. Overall Take-aways}
\begin{itemize}\itemsep0pt
\item \textbf{Opus 4.6} is the best information-recall persona simulator: lowest Modal MAE and EAD on Q2--Q9 across every subgroup.
\item \textbf{GPT-5.2} is the best perception simulator under the EAD criterion and the closest to humans by raw K--S distance.
\item \textbf{Sonnet 4.5} hits modal perception answers correctly but is the least faithful overall on the Information block.
\item \textbf{GPT-4.1} is consistently mid-pack and statistically separated from the leaders on most comparisons.
\item All ten-iteration averages have negligible SEM relative to inter-model gaps, so the ordering above is robust to sampling noise.
\item The \emph{Human (split-half)} baseline reveals that all models still sit well above the human-disagreement floor: the closest model is roughly $1.5\times$--$2\times$ the within-human noise level on K--S distance, leaving clear headroom for future improvements in persona alignment.
\end{itemize}

